\paragraph{Direct joint conditioning of the complete target.}
For each completely matured block, reproject its issued source probabilities
and complete outcomes through the current fixed candidate/reference versions.
Record the forecast coordinates that were actually observed and the complete
normalized target $U_j=g_j/N_j$.  The target is an observed coordinate even
when historical source forecasts are missing.  CJ-T models the joint vector
$X=(H_1,\ldots,H_m,U)$ directly; it requires no independence of $U$ and the
residual $H-U\mathbf1$.

Missing historical forecast coordinates are integrated with a declared
Gaussian Schur law.  If $E_j$ selects the observed coordinates, including $U$,
the completed working moments are
\begin{equation}
 \widetilde x_j=\lambda+
 C E_j^{\mathsf T}(E_j C E_j^{\mathsf T})^{-1}(x_{j,O}-E_j\lambda),\qquad
 V_j=C-C E_j^{\mathsf T}(E_j C E_j^{\mathsf T})^{-1}E_j C.
 \label{cjt_masked_joint_completion}
\end{equation}
The observed forecasts and target are unchanged.  The missing-coordinate
means and uncertainty are working moments, not additional observations.
Retain the original relative and absolute relevance weights.  A common prior
atom and positive definite smoothing give a joint Gaussian mixture
\begin{equation}
 f(h,u)=\sum_{j=0}^{J}\alpha_j
 \phi_{m+1}((h,u);\lambda_j,\Sigma_j),\qquad
 \alpha_j>0,\quad\sum_j\alpha_j=1,\quad\Sigma_j\succ0.
 \label{cjt_joint_target_mixture}
\end{equation}
This is a working density on $\mathbb R^m\times\mathbb R$.
The source likelihood is evaluated at the physical $h\in[-1,1]^m$; its
Gaussian tails outside that source domain are an approximation.
Only the complete target is restricted to $[-1,1]$:
\begin{equation}
 p(u\mid h)=\frac{\mathbf1_{[-1,1]}(u)f(h,u)}
 {\int_{-1}^{1}f(h,v)\,dv}.
 \label{cjt_bounded_direct_target_posterior}
\end{equation}
The target prior is the fitted marginal of $f$, rather than a separately
imposed uniform prior.  This construction does not establish physical
conditional calibration under drift or a correct missing-coordinate law.

\begin{proposition}[Analytic direct-target mixture conditioning]
\label{cjt_analytic_target_conditioning}
Partition each component mean and covariance by $(H,U)$.  Put
\begin{align}
 \eta_j(h)&=\lambda_{j,U}+\Sigma_{j,UH}\Sigma_{j,HH}^{-1}
                 (h-\lambda_{j,H}),\nonumber\\
 v_j&=\Sigma_{j,UU}-\Sigma_{j,UH}\Sigma_{j,HH}^{-1}\Sigma_{j,HU}>0,\nonumber\\
 Z_j(h)&=\Phi((1-\eta_j(h))/\sqrt{v_j})
          -\Phi((-1-\eta_j(h))/\sqrt{v_j}).
 \label{cjt_conditional_component_parameters}
\end{align}
The bounded posterior is a mixture of
$\mathcal N(\eta_j(h),v_j)$ restricted to $[-1,1]$, with weights
\begin{equation}
 \rho_j(h)=\frac{\alpha_j\phi_m(h;\lambda_{j,H},\Sigma_{j,HH})Z_j(h)}
 {\sum_\ell\alpha_\ell\phi_m(h;\lambda_{\ell,H},\Sigma_{\ell,HH})Z_\ell(h)}.
 \label{cjt_conditional_component_weights}
\end{equation}
\end{proposition}
\begin{proof}
Gaussian conditioning factors each component joint density into its forecast
marginal density and its scalar conditional target density.  Integration over
the target bounds supplies $Z_j$, and normalization gives $\rho_j$.
\end{proof}
The component forecast evidence and target truncation mass both matter.
Uniformly averaging conditional components would generally change the law.
The issued mean and scale are $F=N\mathbb E[U\mid h]$ and
$S=N\sqrt{\operatorname{Var}(U\mid h)+\nu^2}$ under the same paid gate.

\paragraph{An exact source--target marginal information control.}
Let $f_U$ be the raw joint target marginal and $f_{sU}$ its source--target
bivariate marginals.  Positive Gaussian smoothing ensures $f_U(u)>0$.
Define
\begin{equation}
 f_{\mathrm P}(h,u)=f_U(u)\prod_{s=1}^m f(H_s=h_s\mid U=u)
 =\frac{\prod_{s=1}^m f_{sU}(h_s,u)}{f_U(u)^{m-1}}.
 \label{cjt_exact_bivariate_product_control}
\end{equation}
\begin{proposition}[Preserved target and source--target information]
\label{cjt_product_preserves_bivariate_information}
The density $f_{\mathrm P}$ integrates to one and has exactly the same $f_U$
and every $f_{sU}$ as $f$.  Common target restriction to $[-1,1]$ preserves
this equality, and either law has the same target normalizing constant.
For one active source the joint and product target posteriors coincide.
\end{proposition}
\begin{proof}
Each $f(H_s\mid U=u)$ integrates to one.  Integrating all forecast coordinates
gives $f_U$, and integrating all but coordinate $s$ gives $f_{sU}$.
If $C=\int_{-1}^{1}f_U(u)\,du$, restricted bivariate marginals are
$\mathbf1_{[-1,1]}(u)f_{sU}/C$, and the restricted target marginal is
$\mathbf1_{[-1,1]}(u)f_U/C$.  On $[-1,1]$, their quotient in
Eq.~\eqref{cjt_exact_bivariate_product_control} is
$f_{\mathrm P}/C$, extended by zero outside the target bounds;
the repeated target-prior factors cancel.
The singleton case is the original bivariate density itself.
\end{proof}
The control keeps each source's nonlinear association with the complete target
and removes only dependence among forecast coordinates conditional on that
target.  Multiplying the bivariate densities without the denominator would
count the target prior $m$ times and would not be this control.
Its bounded target moments use scalar quadrature with refinement checks.

The full-Gaussian control instead uses the exact raw mixture mean and covariance
$\bar\lambda=\sum_j\alpha_j\lambda_j$ and
$\bar\Sigma=\sum_j\alpha_j[\Sigma_j+
(\lambda_j-\bar\lambda)(\lambda_j-\bar\lambda)^{\mathsf T}]$,
then performs the same direct-target conditioning and target restriction.
For a genuinely Gaussian working joint law it coincides with CJ-T.
Any extra paid action supplied by retained joint shape must be evaluated under
the same model versions, complete service accounting and execution ledger.
The fixed-state oracle information gap and fitted paid-policy regret bound
apply to this model as before; neither density retention nor numerical accuracy
proves a bound on the unknown physical conditional mean error.
